%% This file contains the additional packages and custom macros migrated from
%% the original arXiv version (cmu document class) to the CVPR 2026 template.
%% Packages already loaded by cvpr.sty (times, xspace, xcolor[dvipsnames],
%% graphicx, amsmath, amssymb, booktabs, natbib, caption, subcaption, url,
%% color, lineno, cleveref, enumitem) are not re-imported here.

%%
%% Inline annotations
%%
\newcommand{\red}[1]{{\color{red}#1}}
\newcommand{\todo}[1]{{\color{red}#1}}
\newcommand{\TODO}[1]{\textbf{\color{red}[TODO: #1]}}
%% Disabled for blind-review submission: keep TODO/todo silent in case any
%% accidental call survives in the source tree.
\renewcommand{\TODO}[1]{}
\renewcommand{\todo}[1]{#1}

%%
%% Math / typesetting packages
%%
\usepackage{microtype}
\usepackage{mathtools}
\usepackage{amsthm}
\usepackage{mathrsfs}
\usepackage{nicefrac}
\usepackage{dsfont}
\usepackage{bbm}
\usepackage{float}
\usepackage{tabularx}
\usepackage{xr}

%%
%% Tables / figures
%%
\usepackage{multirow}
\usepackage{makecell}
\usepackage{rotating}
\usepackage{arydshln}
\usepackage{colortbl}
\usepackage[export]{adjustbox}
\usepackage{wrapfig}
\usepackage{pifont}
\usepackage{soul}

%%
%% Code / boxes
%%
\usepackage{listings}
\usepackage{tcolorbox}
\tcbuselibrary{skins,breakable}

%%
%% Algorithms
%%
\usepackage{algpseudocode}
\usepackage{setspace}

%%
%% hypcap is intentionally not loaded here; it must come after hyperref and is
%% unnecessary for the CVPR review ruler.

%% --- Dashed lines tuning (arydshln) ---
\setlength{\dashlinedash}{2pt}
\setlength{\dashlinegap}{2pt}

%% --- Colored check / cross symbols and helpers ---
\newcommand{\cmark}{\textcolor{green!60!black}{\ding{51}}}
\newcommand{\xmark}{\textcolor{red!70!black}{\ding{55}}}
\definecolor{softred}{RGB}{200,10,10}
\newcommand{\best}[1]{\textbf{\textcolor{softred}{#1}}}
\newcommand{\pmstd}[1]{\scriptsize{\,\ensuremath{\pm}\,#1}}
\newcommand{\task}[2]{\makecell{\text{{#1}}\\\text{{#2}}}}
\newcommand{\meth}[2]{\makecell[l]{\texttt{#1}\\\texttt{#2}}}
\newcommand{\mone}[1]{\makecell[l]{\texttt{#1}}}
\newcommand{\mema}[2]{\makebox[2cm][l]{{#1}} #2}
\newcommand{\memame}[2]{\makebox[2.8cm][l]{#1}  #2}

\definecolor{rowgray}{gray}{0.96}
\definecolor{cellred}{RGB}{230,30,30}
\newcommand{\bestbest}[1]{\fcolorbox{red}{red!5}{\best{#1}}}
\definecolor{headerblue}{RGB}{215,220,255}

\newcounter{q}
\newcommand{\Q}{Q\stepcounter{q}\arabic{q}}

%% --- Project / task name macros ---
\newcommand{\TSone}{\texttt{Counting}}
\newcommand{\TStwo}{\texttt{Permanence}}
\newcommand{\TSthree}{\texttt{Reference}}
\newcommand{\TSfour}{\texttt{Imitation}}
\newcommand{\benchname}{RoboMME}
\newcommand{\modelname}{MME-VLA}
\newcommand{\trainsteps}{770k}

%% --- Color palette used in figures and listings ---
\definecolor{deepblue}{rgb}{0,0,0.5}
\definecolor{deepred}{rgb}{0.6,0,0}
\definecolor{deepgreen}{rgb}{0,0.5,0}
\definecolor{rliableolive}{HTML}{BBCC33}
\definecolor{rliableblue}{HTML}{77AADD}
\definecolor{rliablered}{HTML}{EE8866}
\definecolor{myblue}{rgb}{0.1,0.4,0.9}
\definecolor{lightblue}{rgb}{0.22,0.45,0.70}
\definecolor{figurebackground}{HTML}{F9F8F3}

\definecolor{highlightmistake}{RGB}{255, 179, 179}
\definecolor{highlightcorrect}{RGB}{179, 255, 179}

%% --- Listings (JSON style) ---
\lstdefinelanguage{json}{
  morestring=[b]",
  stringstyle=\ttfamily,
  showstringspaces=false,
  breaklines=true,
  basicstyle=\ttfamily\footnotesize,
  keywords={true,false,null},
  keywordstyle=\bfseries,
}

\lstset{
  columns=fullflexible,
  keepspaces=true,
  breaklines=true,
  frame=none,
  xleftmargin=0pt,
  aboveskip=0pt,
  belowskip=0pt,
}

%% --- AI box (used for "Key Takeaways" callouts) ---
\tcbset{
  aibox/.style={
    width=\linewidth,
    top=8pt,
    bottom=4pt,
    colback=figurebackground,
    colframe=black,
    colbacktitle=black,
    enhanced,
    center,
    attach boxed title to top left={yshift=-0.1in,xshift=0.15in},
    boxed title style={boxrule=0pt,colframe=white,},
  }
}
\newtcolorbox{AIbox}[2][]{aibox,title=#2,#1}

%% --- Theorem-style environments ---
\theoremstyle{plain}
\newtheorem{theorem}{Theorem}[section]
\newtheorem{proposition}[theorem]{Proposition}
\newtheorem{lemma}[theorem]{Lemma}
\newtheorem{corollary}[theorem]{Corollary}
\theoremstyle{definition}
\newtheorem{definition}[theorem]{Definition}
\newtheorem{assumption}[theorem]{Assumption}
\theoremstyle{remark}
\newtheorem{remark}[theorem]{Remark}

\hyphenation{pre-print}
